\documentclass[reprint,superscriptaddress,nofootinbib,nobibnotes,amsmath,amssymb,aps,prd,floatfix,longbibliography]{revtex4-2}
\usepackage[T1]{fontenc}
\usepackage[utf8]{inputenc}
\usepackage{amsmath,amssymb,bm,booktabs,graphicx,xcolor}
\usepackage{hyperref}
\usepackage{xurl}
\usepackage{placeins}
\def\btemubody{body.tex}
\def\btemubib{references}
\graphicspath{{figures/}}
\hypersetup{colorlinks=true,linkcolor=blue,citecolor=blue,urlcolor=blue}
\providecommand{\tightlist}{\setlength{\itemsep}{0pt}\setlength{\parskip}{0pt}}
\setlength{\emergencystretch}{2em}
\renewcommand{\dbltopfraction}{0.95}
\renewcommand{\textfraction}{0.02}
\begin{document}
\title{Small-scale matter spectra and reionization: sensitivity and galaxy freedom}
\author{Haina Huang}
\affiliation{Department of Physics, The Chinese University of Hong Kong, Shatin, Hong Kong, China}
\date{November 6, 2026}
\begin{abstract}
Small-scale matter spectra and galaxy formation both affect the ionizing sources of reionization. We compare fixed-galaxy histories and independently fitted luminosity functions (LFs) in a phenomenological spectrum model with break scale $k_p$ and slope $m_s$ at fixed $\sigma_8$. The exported histories separate as the spectrum changes. At larger break scales, the LF archive contains candidate fits close to the power-law (PL) reference. Under the adopted galaxy prescription, the released six-coordinate samples at $k_p=1\,h\,\mathrm{Mpc}^{-1}$ and $m_s=1.5$ give lower active-halo star-formation factors than PL, with BT-to-PL ratios of marginal median star-formation factors 0.123, 0.195 and 0.308 at halo masses $10^9,10^{10},10^{11}M_\odot$. These calculations establish conditional history sensitivity and fit compatibility; separating weak direct sensitivity from galaxy compensation requires a matched observable comparison.
\end{abstract}
\maketitle
\begin{quote}\small Physics working draft v0.22.\end{quote}
\input{\btemubody}
\end{document}
